

\section{Related Work}
\label{Related Work}


\subsection{Test-Time Adaptation}

Test-Time Adaptation (TTA) aims to mitigate distribution shifts by updating a pretrained model using unlabeled test data without accessing source-domain samples. With the rapid development of multimodal learning, recent studies have extended TTA from single-modality settings to multimodal scenarios, where heterogeneous noise patterns and dynamic modality interactions introduce additional challenges. Early approaches mainly relied on entropy minimization or normalization-layer updates for efficient online adaptation. More recent works explore multimodal TTA by incorporating cross-modal attention modulation, modality reliability estimation, and feature alignment strategies to improve robustness under noisy environments. 

In real-world applications, however, target domains often evolve continuously and exhibit non-stationary distributions. To address this issue, Continual Test-Time Adaptation (CTTA) has been proposed to enable stable long-term adaptation under streaming target data. Representative methods introduce mechanisms such as teacher--student consistency, dynamic adapters, or expert-based structures to maintain stability across changing domains. Meanwhile, prompt-based adaptation has emerged as a parameter-efficient paradigm, where lightweight prompts are optimized instead of updating the full model parameters. Forward-only optimization and visual prompting strategies further improve efficiency, making prompt-based methods particularly suitable for resource-constrained multimodal deployment scenarios.

\subsection{Test-Time Adaptation for Multi-Modal Data}

Although significant progress has been made in multimodal TTA and CTTA, existing studies mainly focus on improving cross-modal robustness through attention reweighting, modality selection, or statistical calibration. For instance, reliability-aware methods dynamically suppress noisy modalities, while alignment-based approaches enhance fusion stability by matching feature distributions across modalities. In addition, several CTTA frameworks introduce teacher--student structures, expert fusion, or memory replay mechanisms to alleviate error accumulation during continual adaptation. These strategies improve robustness under complex noise and missing-modality scenarios, but most of them rely on local adjustments or short-term statistics, which may limit their ability to capture long-term cross-modal co-evolution in non-stationary environments.

However, current methods still struggle to balance stability and efficiency. Many multimodal TTA approaches require explicit gradient updates or complex architectural designs, while existing prompt-based methods often emphasize single-step statistical alignment without explicitly modeling dynamic cross-modal consistency. As a result, they may suffer from semantic drift or unstable fusion during long-term adaptation. To address these limitations, we propose a calibrated prompt-based continual adaptation framework for non-stationary multimodal environments, which introduces dynamic consistency calibration to achieve robust and efficient adaptation.

